\chapter{总结与展望}

\section{本文研究工作总结}

本文围绕多主体信息传播影响力评估这一总目标，从消息层、路径层和交互类型层三个评估维度展开研究，设计并实现了多主体信息传播影响力评估系统。主要工作总结如下：

（1）构建了包含主体类型标注的社交媒体传播数据集。基于X平台，通过"热点话题获取$\to$根推文检索$\to$评论树采集"的三阶段流水线，构建了包含人类用户与AI代理（Grok）评论的多主体数据集。经去重与过滤后，用于传播价值评估的数据集包含158,972条评论（Human 155,529条、Grok 3,443条），用于传播延续预测的数据集包含166,543个节点（Human 163,055个、Agent 3,488个）。

（2）提出了基于组感知融合的评论传播价值评估模型。针对消息层评估维度，提出Group-aware Fusion TabTransformer模型，将发布环境、文本语义、情感倾向和主体标识四组异构特征通过门控融合机制进行样本级自适应加权，经Transformer编码器捕捉组间高阶交互后输出传播价值评分。模型在三项指标上均优于四个基线；消融实验表明，去除主体标识特征后Agent子集Recall@P=0.5骤降19.0\%，验证了显式引入主体类型特征的必要性。

（3）提出了基于交互感知的异构路径传播延续预测模型。针对路径层与交互类型层评估维度，提出交互感知的异构路径编码器（IHPE），以根节点到目标节点的完整路径为输入，通过交互类型编码层显式建模四种主体交互模式（H$\to$H / H$\to$A / A$\to$H / A$\to$A），经LSTM捕捉路径上下文依赖后预测2跳增长概率。模型在三项指标上均取得最优；消融实验表明，交互类型编码是核心贡献组件，移除后R@P0.8骤降36.4\%。

（4）设计并实现了多主体信息传播影响力评估系统。系统采用四层架构与六个功能模块，实现了从数据导入到评估结果可视化的端到端流程。多主体对比分析模块将两个模型的输出按主体类型聚合，从传播价值、传播延续和情感倾向三个维度量化两类主体的传播影响力差异，并生成结构化的自然语言分析报告，实现了"从多主体视角系统性量化信息传播影响力差异"的总目标。

\section{未来工作展望}

本文在多主体信息传播影响力评估方面进行了初步探索，但受限于数据规模、平台范围和方法能力，仍存在若干可改进之处，未来可从以下方向展开深入研究：

（1）扩展数据来源与主体类型覆盖。本文的数据集仅来源于X平台，AI代理类型限于Grok单一主体。未来可将数据采集扩展至Reddit、YouTube等多个社交媒体平台，纳入更多类型的AI代理（如ChatGPT、Claude等大语言模型驱动的社交账号），构建跨平台、多AI代理类型的数据集，以验证模型在不同平台环境和多种AI代理共存场景下的泛化能力。同时，可探索自动化的主体类型识别方法，降低对人工标注的依赖。

（2）引入动态时序建模与在线预测能力。本文的两个评估模型均基于评论发布时刻的静态特征进行预测，未利用评论发布后的动态反馈信息（如点赞数随时间的变化、早期回复的到达模式等）。未来可引入动态时序建模方法，将早期传播反馈作为增量输入，实现随时间推移不断更新的在线预测，从而提高预测精度并支持实时传播监测场景。

（3）融合图神经网络进行评论树全局建模。本文的传播延续预测模型以单条路径为输入单元，未同时利用评论树中其他分支的信息。实际上，同一评论树内不同讨论分支之间可能存在话题关联和竞争效应，某条分支的活跃度会影响其他分支的传播延续。未来可探索基于图神经网络的评论树全局建模方法，在保留路径级异构交互信息的同时，捕捉跨分支的结构依赖关系，进一步提升传播延续预测的准确性。